\section*{Works cited and source verification}
\addcontentsline{toc}{section}{Works cited and source verification}
\begin{thebibliography}{9}
\bibitem{sprague1935}
R. Sprague, ``\"Uber mathematische Kampfspiele,'' \emph{T\^ohoku Mathematical Journal}, first series, vol.~41 (1935), pp.~438--444.\newline
Original journal record: \url{https://www.jstage.jst.go.jp/article/tmj1911/41/0/41_0_438/_article}. The metadata and page span were checked against the journal listing.
\bibitem{grundy1939}
P.~M. Grundy, ``Mathematics and Games,'' \emph{Eureka}, no.~2 (1939), pp.~6--8.\newline
Bibliographical entry and link to scanned source at \url{https://oeis.org/A002188/a002188.pdf}; the bibliographical description is corroborated by OEIS. The scan was not used for any copied proof.
\bibitem{winningways2001}
E.~R. Berlekamp, J.~H. Conway, and R.~K. Guy, \emph{Winning Ways for Your Mathematical Plays}, vol.~1, second ed., A~K Peters, 2001.\newline
Publisher record: \url{https://www.routledge.com/Winning-Ways-for-Your-Mathematical-Plays-Volume-1-2nd-Edition/Berlekamp-Conway-Guy/p/book/9781568811307}; verified against publisher and author's publication listings.
\bibitem{gurvich2012}
V.~Gurvich, ``On Repetition-free Subtraction Games and Vile-dopey Integers,'' RUTCOR Research Report RRR-01-2012, Rutgers University, January 2012.\newline
Author's publication list: \url{https://rutcor.rutgers.edu/~gurvich/publist.pdf} (entry 240). A report transcription was located, but its detailed theorems were \emph{not} invoked here.
\end{thebibliography}
\textit{Source-use boundary.} The P/N and Sprague--Grundy methods are standard background and are credited above. All specific interval, finite-window, template-obstruction and finite-range arguments in these documents are written self-contained from the game definition, not borrowed from an external proof. Originality relative to \emph{all} prior literature has not been established; a thorough review remains necessary.
